\chapter{Front and back buffers}
\label{ch:buffers}

\begin{aside}
Double-buffering is older than computers. A scribe writing on
parchment kept a wax tablet for drafts; the parchment only saw the
final version. \maya{} does the same with terminal cells: drafts
land on the back buffer, the final version flushes to the screen.
\end{aside}

The canvas in Chapter~\ref{ch:canvas} introduced packed cells. This
chapter looks at why \maya{} keeps \emph{two} canvases instead of
one, and what each is for.

\section{The problem with one buffer}

Imagine the renderer only had a single grid. To redraw the screen
after a state change, it would have to:

\begin{enumerate}[leftmargin=*]
  \item Clear the grid.
  \item Walk the element tree and paint into it.
  \item Emit ANSI for every cell, because there is nothing to
    diff against.
\end{enumerate}

Step three is the killer. A modest 200x60 terminal has 12,000
cells. At roughly 20 bytes of ANSI per cell (cursor move, SGR set,
glyph, reset) that is 240 KB of output per frame. At 60 frames per
second that is 14 MB/s of escape codes flowing through the pty.
Most terminals choke long before that.

\section{Two buffers, one diff}

With two buffers, the renderer paints into the \emph{back} buffer
and compares it to the \emph{front}. Only cells that differ become
ANSI output. A typing animation in a text field touches one cell;
the resulting frame is roughly 30 bytes of ANSI, not 240 KB.

\begin{lstlisting}
class Renderer {
    Canvas front_;   // what the screen currently shows
    Canvas back_;    // what we want it to show
    // ...
};
\end{lstlisting}

The diff loop is the subject of Chapter~\ref{ch:simd}. Here we focus on
the \emph{lifecycle} of the two buffers around a frame.

\section{Frame lifecycle}

A frame in \maya{} proceeds in five steps:

\begin{enumerate}[leftmargin=*]
  \item \textbf{Resize check.} If the terminal was resized, both
    buffers are reallocated to the new size. The front is filled
    with a sentinel cell that cannot equal any real cell, forcing
    the next diff to redraw everything.
  \item \textbf{Clear back.} The back buffer is filled with the
    default cell (space, default fg/bg).
  \item \textbf{Paint.} The element tree is walked depth-first and
    each element writes into the back buffer.
  \item \textbf{Diff and emit.} The two buffers are compared; for
    each differing run, ANSI is emitted.
  \item \textbf{Swap.} The back becomes the new front. The next
    frame begins.
\end{enumerate}

The swap is logical, not physical: \maya{} does not memcpy. It
swaps pointers (or, equivalently, swaps the role of two storage
slots). Memcpy of 12,000 cells $\times$ 8 bytes = 96 KB per swap
would dwarf the diff savings.

\section{Sentinel cells}

\begin{idea}
\textbf{Sentinel.} A value that cannot occur in normal use, used
to mark a slot as ``not yet computed'' or ``needs full redraw''.
\end{idea}

After resize, \maya{} fills the front buffer with a sentinel cell
whose width field is set to 3 --- impossible for any real cell
(widths are 0, 1, or 2). The diff sees every back cell as different
and redraws the whole screen.

The same trick handles the very first frame: the constructor
initialises the front to sentinels, so the first diff is a full
paint.

\section{Why not triple-buffer?}

Some GUI frameworks triple-buffer: while one buffer is on screen,
another is being painted, and a third is being diffed. The aim is
to never block the paint thread on the screen.

\maya{} does not need this. Painting and emitting happen on the
same thread, and the terminal is the slow link, not the CPU. Two
buffers are enough.

\begin{funfact}
Triple-buffering shows up in OpenGL as \code{GL\_TRIPLE\_BUFFER}
and in Wayland as the \code{wl\_surface} commit protocol. The
terminal world skipped straight from no buffering (telexes,
printing terminals) to double-buffering (curses, 1980).
\end{funfact}

\section{Cell-level write barriers}

Painting into the back buffer is a write. \maya{}'s painter
exposes a small API:

\begin{lstlisting}
void put(int x, int y, Cell c);
void put_str(int x, int y, std::string_view s, Style st);
void fill(Rect r, Cell c);
\end{lstlisting}

There are no random writes outside this API. The reason is
clipping: every \code{put} checks bounds against the current
clip rectangle. If an element draws outside its allotted space
the write is silently dropped. Without this, a misbehaving
element could corrupt its neighbours.

\begin{warning}
The clip rectangle is per-element, set by the layout engine
before paint. If you ever write your own paint method (custom
widget), you must respect it. \maya{}'s built-in elements do
this for you.
\end{warning}

\section{Damage tracking}

The diff in Chapter~\ref{ch:simd} is the simplest form of damage
tracking: compare every cell. Other libraries (notch, blessed)
maintain explicit damage rectangles --- ``cells x1..x2, y1..y2
changed''.

\maya{} chose full-grid diff because:

\begin{itemize}[leftmargin=*]
  \item It is simpler. No bookkeeping at write time.
  \item It is SIMD-friendly: the inner loop is a wide
    \code{vpcmpeqq} over 4 or 8 cells at a time.
  \item For typical TUI sizes (under 30,000 cells) the diff is
    sub-millisecond on any laptop made since 2015.
\end{itemize}

The trade-off is that explicit damage rectangles let you skip
\emph{painting} entire regions, not just emitting. \maya{} pays
the paint cost every frame and trusts the diff to keep the wire
quiet.

\section{Buffer ownership}

Both buffers are owned by the \code{Renderer}. The element tree
borrows a painter handle for the duration of a frame but never
sees the canvas storage directly.

\begin{lstlisting}
class Painter {
    Canvas& back_;
    Rect    clip_;
public:
    void put(int x, int y, Cell c);
    // ...
};
\end{lstlisting}

This makes ownership obvious and makes the painter cheap to
construct (one reference, one rect). Elements never need to
allocate during paint.

\section{Allocation strategy}

A canvas is a flat \code{std::vector<uint64\_t>} of size
\code{width * height}. Resizes use \code{vector::assign} which
amortises to one allocation per resize burst. Steady-state has
zero allocations during paint or diff.

\begin{tryit}
Run any \maya{} program and resize the terminal aggressively
with the mouse. The redraw is smooth because each resize is
one allocation and one full-screen redraw, not a flood of
per-cell operations.
\end{tryit}

\section{What you do not buffer}

The cursor position is not buffered. Neither is the title bar.
Neither are bell, hyperlinks (the URL string itself; the slot
in the cell is), or mouse cursor mode.

These are emitted directly at the end of each frame:

\begin{lstlisting}
emit("\x1b]0;" + new_title + "\x07");  // OSC 0
emit(cursor_move(cx, cy));
\end{lstlisting}

The reason: there is only one of each, and they are tiny. A
buffer for state with one slot is just a variable.

\section{Recap}

\begin{recap}
\begin{itemize}[leftmargin=*]
  \item Two buffers let \maya{} send only changes, not full
    frames.
  \item The back is the draft, the front is the screen.
  \item Resize seeds the front with sentinels so the next diff
    redraws everything.
  \item The clip rectangle is the only write barrier; elements
    cannot scribble outside their box.
  \item Cursor and title are unbuffered: one slot needs no
    buffer.
\end{itemize}
\end{recap}

\section*{Workshop}

\begin{enumerate}[leftmargin=*]
  \item Add a counter to your app that prints
    \code{frame N: K cells changed} on every frame. K is the
    output of the diff. Type into a text field and watch K
    stay small.
  \item Force a full redraw by writing the sentinel into the
    front buffer (you will need to expose a debug method).
    Confirm the next frame's K equals the screen area.
  \item Time the paint phase (back buffer write) and the
    diff phase separately. Which is dominant in your app?
\end{enumerate}
